\chapter*{Glossary}
\addcontentsline{toc}{chapter}{Glossary}
\markboth{Glossary}{}

This book is written for two trades at once, so a term that is obvious to
one reader is often new to the other. Every entry here is the book's own
usage, taken from the chapter that introduces it, and each term appears in
exactly one form throughout: where two words are commonly treated as
synonyms and this book does not, the entry says so. The numbers in the
entries restate figures the chapters give, and each is sourced where its
chapter gives it, not here.

% A5 measure is narrow and many entries carry long hyphenated compounds and
% parenthesised acronyms. Give TeX room to stretch rather than overrun.
\begingroup
\emergencystretch=3em
\hbadness=6000
\begin{description}[leftmargin=0pt,style=unboxed,itemsep=0.35em,parsep=0.15em]
\item[12CP (Twelve Coincident Peak)] \hypertarget{glossary:12cp}{}The commission's proposed replacement for 4CP: twelve monthly peaks, one for every month from October through September, each measured over a thirty-minute interval instead of fifteen. \emph{Proposed in the Texas Register on July 24, 2026, final rules required by December 31, 2026. Spreading the same annual charge over twelve peaks makes each worth about a third as much, roughly \$3,500 per megawatt at the Oncor floor instead of \$10,500. See 4CP, minimum billing demand.}
\item[2x16] \hypertarget{glossary:2x16}{}The sixteen daytime hours, 6 a.m. to 10 p.m. Central time, on Saturdays, Sundays and the six market holidays: every daytime hour that on-peak leaves out. \emph{The third of the three standard block products, with on-peak (5x16) and off-peak nights (7x8), and 1,760 hours in 2025. At LZ\_WEST that year the window averaged \$32 per megawatt-hour, the cheapest of the three, because it combines the solar belly with weak weekend demand. See block.}
\item[4CP (Four Coincident Peak)] \hypertarget{glossary:4cp}{}How ERCOT allocates transmission costs. In each summer month, June through September, ERCOT finds the single 15-minute interval of highest system-wide demand; a large consumer's share of the next year's transmission bill is then set by its average draw across exactly those four intervals. \emph{Demand here is a settlement quantity, not the raw control-room load: ERCOT sums metered generator output, nets off DC-tie flows and excludes grid-scale battery charging. It confirms each interval only after the month ends, so loads stay off across several candidate afternoons. Oncor's own share was about \$21,400 per megawatt of 4CP demand a year in 2025, and the ERCOT-wide charge a transmission-voltage site bears is roughly two to three times that, depending on the wires utility. The charge falls on peak megawatts, not on energy, so what avoidance is worth per megawatt-hour depends on the site: Chapter~6's ledger puts it at about \$5.65/MWh of the delivered price, at Oncor's transmission rate weighted over 2025 and a 95 percent capacity factor. Distinct from the summer window.}
\item[aggregator] \hypertarget{glossary:aggregator}{}A company that bundles many customers' curtailable load and fronts them to ERCOT through its own QSE, handling telemetry, offers, tests and settlement, and keeping a share of the proceeds. \emph{What it charges is a private commercial term: no aggregator publishes its revenue share, so this book prints none. The route costs a slice of the revenue and surrenders control of when the site is offered.}
\item[ancillary service (AS)] \hypertarget{glossary:ancillary-service}{}Insurance ERCOT buys alongside energy: standing capabilities that keep the lights on when something breaks. ERCOT's current set has five members; a sixth, the Dispatchable Reliability Reserve Service, is approved but not yet running, and only generators may sell it. \emph{Because ERCOT buys readiness, a load that can vanish in seconds is itself a sellable service. Since RTC+B every SCED run awards and prices all five; the Day-Ahead Market still clears them, now as financial positions.}
\item[Antminer S19j Pro (S19j Pro)] \hypertarget{glossary:antminer-s19j-pro}{}One of the book's two standing machines: a 2021-vintage workhorse that still populates mining farms in large numbers, 104~TH/s at 3,068 W, which is 29.5~J/TH. \emph{At the book's \$42.18 price anchor its breakeven hashprice is \$0.030 per terahash-day, a line the market crossed in 2026. At the 2025 mean hashprice its strike was \$72 and it paused 14 percent of the year: the machine that earns its keep by knowing when to stop.}
\item[Antminer S21 XP (S21 XP)] \hypertarget{glossary:antminer-s21-xp}{}The other of the book's two standing machines: the 2024 flagship of the same manufacturer, 270~TH/s at 3,645 W, which is 13.5~J/TH. \emph{At the same price anchor its breakeven hashprice is \$0.014 per terahash-day, below anything the market has recorded; at the 2025 mean hashprice its \$157 strike sat above almost every price the year produced. The 3,645 W is wall draw, not the PSU rating. It inherits the S19j~Pro's 2025 strike when hashprice falls to 46 percent of its 2025 level.}
\item[ASIC] \hypertarget{glossary:asic}{}An Application-Specific Integrated Circuit: a chip, and by extension the machine built around it, engineered for nothing but mining, capable of immense numbers of hashes per second at the cost of significant energy consumption. \emph{Replaced general-purpose computers in mining entirely by 2015. An ASIC cannot be repurposed for other computing work, which is why an idled mining machine is a load that can stop, not a data center that can switch tasks.}
\item[autotuning firmware] \hypertarget{glossary:autotuning-firmware}{}Miner firmware that retunes each hashing chip's voltage and frequency individually, so the machine can be held at any power level across a wide range rather than at the one operating point stock firmware ships with. The component doing the work is the autotuner. \emph{Per-chip retuning is why the autotuned fleet average beats the nameplate at stock power, 28.5 against 29.5~J/TH on the S19j~Pro. Moving to a power target with no saved tuning profile has a cost: the autotuner walks every chip to new settings, and the machine hashes below its optimum for the settling period of that tuning run. A target it already holds a profile for can be applied again without a new run, but only that exact target. After its first move away from stock power, the loop in Chapter~5's backtests applies only four targets on the S21~XP and eight on the S19j~Pro, so each needs one tuning run, barring a re-tune after a large change in chip temperature.}
\item[availability factor] \hypertarget{glossary:availability-factor}{}The share of an ERS load's committed quarter-hours in which it was actually available, which sets the availability component of its payment. A load consuming less than 95 percent of the capacity it contracted is counted unavailable for that interval, whatever the reason. \emph{Real deployments and tests are carved out; price response and peak-dodging are not. Only when a QSE's whole portfolio falls below 95 percent for the contract period can ERCOT suspend it and refer it to the commission. A summer afternoon block runs 255 obligated hours, so the 5 percent slack is under 13 hours.}
\item[backstop] \hypertarget{glossary:backstop}{}The rule that owns pause and resume. It reacts to accumulated damage rather than to a price merely crossing breakeven: while the price sits above the machine's breakeven at its most efficient power target within the loop's operating range the excess integrates, and the fleet pauses the moment the accumulated excess reaches \$10 per megawatt of load. The integral resets when the price dips back below breakeven, so a shallow blip never trips it, while a single five-minute print \$120 or more above breakeven does. \emph{Resuming is deliberately harder than pausing: the price must fall 10 percent below breakeven and hold 15 minutes, with minimum runtimes of 30 minutes each way and at most eight cycles a day. In thirty months of backtest across all eight ERCOT load zones it tripped 0.12 to 0.53 times a day. Modulation can never trigger a pause.}
\item[balancing authority] \hypertarget{glossary:balancing-authority}{}The grid operator responsible for watching the supply and demand balance in its area and adjusting generation output or shifting loads to keep the two in sync. \emph{The book's own term is grid operator; Chapter~2 gives balancing authority once, as the industry's name for the same role.}
\item[base point] \hypertarget{glossary:base-point}{}A target consumption level sent to a Controllable Load Resource by the same five-minute economic dispatch that steers every power plant in Texas, scored for compliance every interval. \emph{Not the site's own power target. The base point is the grid's instruction; the power target is how the site meets it.}
\item[basis] \hypertarget{glossary:basis}{}The gap between the trading hub price at which a futures block settles and the load zone price a mine actually pays. It is a risk the block does not cover. \emph{Not cost basis: this is a locational spread, not a purchase price. See hub, load zone, block.}
\item[behind the meter] \hypertarget{glossary:behind-the-meter}{}The trade's term for what ERCOT calls a Private Use Network: a mine built beside a generator and taking power across the fence, so that it stops being a retail customer of the grid. \emph{See Private Use Network for the definition and fence for what the deal pays and gives up.}
\item[Bitcoin block] \hypertarget{glossary:bitcoin-block}{}A batch of transactions proposed for addition to the Bitcoin blockchain. Miners repeatedly hash its data in search of a valid proof of work. Not the electricity contract called a block.
\item[black start] \hypertarget{glossary:black-start}{}Restarting a collapsed grid from a handful of designated units that can start with no outside power, typically hydroelectric plants or gas turbines with their own on-site generators. Most plants cannot start at all when the grid is dead, because they need grid electricity to run their own equipment. \emph{Operators use those units to energize pre-planned cranking paths, transmission corridors that carry station power to larger plants, which then start up, pick up a small block of customers and extend the electrified island outward; reconnect too much at once and cold-load pickup, the surge when every thermostat and compressor in a cold, dark city switches on together, collapses the young island's frequency again. This is why recovery takes days, not hours: ERCOT's leadership estimated a full black start in 2021 would have left most of the state dark for weeks.}
\item[block] \hypertarget{glossary:block}{}A fixed price bought for a fixed slice of the week, hedging some hours and leaving the rest exposed. Blocks have exact hour definitions because they trade as monthly futures, each settling against the average of the hub's settlement point prices inside its own window. \emph{Not a Bitcoin block: this is a power product, and the protocol's sense of the word is under block reward, block subsidy and block time. The three standard shapes tile the year with no gap and no overlap: on-peak (5x16), 2x16 and off-peak nights (7x8) were 4,080, 1,760 and 2,920 hours in 2025. A fourth quote, 7x24, is all of them at once. See basis.}
\item[Block Load Transfer (BLT)] \hypertarget{glossary:block-load-transfer}{}A group of customers disconnected from its usual grid area and connected to a neighboring one. \emph{The official peak calculation accounts for these transfers; the book's generation-based demand proxy does not separately reconstruct them.}
\item[block reward] \hypertarget{glossary:block-reward}{}What the miner who finds a valid block receives: newly minted bitcoins plus any transaction fees from the block it processed. This is the revenue that pays for the hardware and the electricity. \emph{The newly minted part alone is the block subsidy; the block reward is subsidy plus fees, and the book never uses one for the other.}
\item[block subsidy] \hypertarget{glossary:block-subsidy}{}The newly created bitcoin paid to the miner of each block, currently 3.125 bitcoin. With 144 blocks in an expected day that is 450 BTC a day paid out to the whole network, before transaction fees. \emph{The larger part of miner revenue, not the whole of it. It is not a price paid by any customer, and no miner can produce more of it.}
\item[block time] \hypertarget{glossary:block-time}{}The average interval between blocks, which the difficulty adjustment holds at roughly 10 minutes no matter how much hashrate the network has. \emph{The 10-minute target is the reference point for every difficulty claim in the book: during Winter Storm Fern average block times stretched to about 12.4 minutes.}
\item[breakeven price] \hypertarget{glossary:breakeven-price}{}The electricity price at which a machine's mining revenue exactly covers its power bill. Above it the machine mines at a loss and, sooner or later, shuts down. \emph{It is the mirror of efficiency at a given hashprice: when the block subsidy alone paid about \$0.032 per terahash-day, as in August 2026, a machine paying \$50 per megawatt-hour breaks even at about 27~J/TH and one paying \$40 at about 33~J/TH. Both move with hashprice. Chapter~4 derives the arithmetic; the same number used as a live tripwire is the strike price.}
\item[capacity market] \hypertarget{glossary:capacity-market}{}A second revenue stream, used in most U.S. markets but not in ERCOT, that pays a generator for simply existing and promising to be available years into the future. \emph{ERCOT has none, and that absence is what forces its scarcity prices to run so high.}
\item[capacity price] \hypertarget{glossary:capacity-price}{}The price of standby itself: dollars per megawatt for each hour an ancillary award is held, written \$/MW per hour, not dollars per megawatt-hour of energy. \emph{A resource earns it on its committed megawatts every hour it holds the award, on top of whatever it earns selling or consuming energy. Keep the capacity-price unit separate from the energy-price unit.}
\item[capture] \hypertarget{glossary:capture}{}The share of a perfect-foresight benchmark's gain that a real strategy earns over the same intervals, the gain being profit above running at stock power and never pausing. \emph{A percentage scored by a strategy, not a price: distinct from captured price, which is the dollars per megawatt-hour a generator's output actually fetched at its node. The shortfall is not all blindness: in Chapter~5's backtests, out of sample, the loop's own rules cost two-fifths to half of it. See oracle (perfect-foresight twin), captured price.}
\item[captured price] \hypertarget{glossary:captured-price}{}The average price a plant's output actually fetched at its own node. \emph{Across ERCOT in 2025 that was \$19 per megawatt-hour for wind and \$21 for solar, less than half the zone price a load pays. It is the floor of the behind-the-meter negotiation: below about \$19 the generator would rather sell to the grid, and above about \$40, priced on the same wind hours, the mine would rather buy those hours from the grid.}
\item[cascading failure] \hypertarget{glossary:cascading-failure}{}A chain reaction in which power from a failed line shifts onto neighboring lines, which overload and trip in turn, until generators and loads across a region fall out of sync. \emph{The 2003 Northeast blackout is the book's case for it.}
\item[the commission] \hypertarget{glossary:the-commission}{}The Public Utility Commission of Texas, the state regulator that sets the rules ERCOT operates under, approves its protocol changes and decides the transmission tariffs this book prices. \emph{It is the body behind the 4CP rule, its proposed 12CP replacement, and the Emergency Response Service. It also selects the market monitor and supervises it, while keeping every enforcement power in its own hands.}
\item[the companion repository] \hypertarget{glossary:the-companion-repository}{}The public code and data behind every number the book computes: each chapter's ledger scripts, the checks on them, and a historical calculator that runs the same engine for a different mix of the two machines, zone or contract. It is at \repourl.
\item[congestion cost] \hypertarget{glossary:congestion-cost}{}The second component of the LMP: what the transmission limits between a node and the rest of the grid add to or subtract from the System Lambda. \emph{When the system is uncongested prices are nearly uniform; when a line reaches its limit the grid splits into cheap and expensive regions and the LMPs say so. In ERCOT this is the only thing location does to a price, because losses are not priced into the LMP.}
\item[Controllable Load Resource (CLR)] \hypertarget{glossary:controllable-load-resource}{}A registration for loads that give ERCOT's dispatch engine direct control in exchange for fuller market participation. Some miners hold it. \emph{A 2024 rule change, approved but awaiting implementation, will move every CLR that is not an aggregation of small sites from zonal to nodal dispatch and settlement. A CLR is a load resource, the controllable form: do not confuse it with the basic form, which trips on an under-frequency relay or drops on a call. A CLR takes a base point, must respond to frequency at five percent droop (under a change ERCOT schedules for early 2027, only a CLR qualified for Regulation or Responsive Reserve), and is the only load registration that can sell Regulation.}
\item[curtailment] \hypertarget{glossary:curtailment}{}Output or load coming off the grid, whoever sets it off. For a load: the site acting on a price or its own rule, or ERCOT or a generator sending an instruction the site has agreed to follow. For a generator it is the mirror image, output held back because the grid cannot take it. This is the industry's word and the book uses it both ways. \emph{Where it matters that the machine stopped on its owner's rule rather than on somebody else's instruction, the book says pause. See pause, load shedding.}
\item[DART spread] \hypertarget{glossary:dart-spread}{}The difference between the day-ahead and real-time prices at the same place and hour. It is a market in itself: buy day-ahead and sell real-time, or the reverse. \emph{In 2024 the two markets differed by about \$17 per megawatt-hour on average in absolute terms.}
\item[Day-Ahead Market (DAM)] \hypertarget{glossary:day-ahead-market}{}A voluntary but financially binding forward market. Participants submit offers and bids for each hour of the following day, ERCOT clears energy and ancillary services together, and the hourly prices are published by 1:30 p.m. Central Time. \emph{A load pays the day-ahead price only for the volumes it actually bought there; whatever it consumes above or below that position settles at the real-time price.}
\item[DC tie] \hypertarget{glossary:dc-tie}{}The direct-current link over which ERCOT exchanges power with neighboring grids. \emph{Power flowing out over the ties is excluded from the settlement demand figure that sets 4CP, one reason the settlement peak and the live dashboard peak can name different quarter-hours; grid-battery charging is the other exclusion and the larger one. ERCOT's whole external connection is three DC ties and one variable-frequency transformer, about 1,220~MW.}
\item[demand response] \hypertarget{glossary:demand-response}{}Grid programs that pay an electricity consumer to stop consuming when the grid is short. A miner enrolled in one earns money from the grid instead of from mining during a price spike. \emph{Demand response is the program that pays for the curtailment, not the act itself. A machine the program's call turns off is curtailed; one its owner's own rule turns off is paused, the self-initiated kind of curtailment. See curtailment, pause.}
\item[demand-side reserve] \hypertarget{glossary:demand-side-reserve}{}A load that stands ready to stop consuming so that its unused power serves the rest of the system, giving the operator the same balancing effect as a reserve generator. A grid operator gets that effect by curtailing a batch of miners that has signed up for it. \emph{The book's one analogy for this is the demand-side battery, an analogy and not a second term, and an imperfect one, since a mine stores nothing and can only stop drawing.}
\item[development fee (devfee)] \hypertarget{glossary:development-fee}{}What Braiins~OS charges for its use, paid in hashrate rather than money: the firmware assigns a share of each machine's hashrate to Braiins Pool, the company's mining pool. \emph{2.5 percent on the two machines of Chapter~9. The ledger does not take it out, and it falls on all of a machine's hashing, not on the power-target layer alone.}
\item[difficulty] \hypertarget{glossary:difficulty}{}The protocol's self-balancing knob. The network automatically adjusts mining difficulty approximately every two weeks, every 2,016 blocks, so blocks keep arriving at the target rate: more hashrate raises difficulty, less hashrate lowers it. \emph{Hashrate and difficulty track each other almost perfectly, because added hashrate directly triggers added difficulty. Difficulty is a consequence of hashrate, not an independent input.}
\item[difficulty adjustment] \hypertarget{glossary:difficulty-adjustment}{}The protocol's two-week recalibration, which resets difficulty to whatever hashrate the network averaged over the previous 2,016 blocks and so restores the ten-minute block cadence at the next boundary. \emph{Winter Storm Fern was the test at scale: about 200~EH/s went offline at U.S.-facing pools, average block times stretched to about 12.4 minutes, the next adjustment was the largest downward move since China's 2021 mining ban, and the one after it rose 14.7 percent as the machines returned. Issuance is safe across it because coins are issued per block, not per minute. Sometimes called the retarget elsewhere; this book uses difficulty adjustment throughout.}
\item[Dispatchable Reliability Reserve Service] \hypertarget{glossary:dispatchable-reliability-reserve-service}{}A sixth ancillary service, approved in July 2026 and waiting on ERCOT's systems: generator capacity that must reach its level within two hours and hold it for four. \emph{Generators only, so a mine's load cannot sell it; only a plant the mine registers as a generation resource could. Its arrival cuts Non-Spin's endurance to two hours.}
\item[distribution] \hypertarget{glossary:distribution}{}The local delivery of power at lower voltages, after transformers step the voltage down, to neighborhoods, businesses and individual consumers. \emph{Large mining sites usually sit on the transmission or sub-transmission system, not on distribution.}
\item[droop] \hypertarget{glossary:droop}{}The standard slope that ties the size of a resource's frequency response to the size of the frequency error; ERCOT sets it at five percent wherever a Controllable Load Resource must answer frequency. \emph{Droop sets how hard a resource answers; the frequency deadband sets where it starts answering. See frequency deadband, Controllable Load Resource.}
\item[duck curve] \hypertarget{glossary:duck-curve}{}The shape of average price plotted against hour of day in a solar-heavy market: a deep belly at midday where solar floods the market and a steep neck in the evening when it withdraws. \emph{At LZ\_WEST in 2025 the midday hours averaged about \$21 and the evening ramp about \$72, more than three times as much a few hours apart on the same average day.}
\item[efficiency curve] \hypertarget{glossary:efficiency-curve}{}A machine's efficiency plotted across its whole power range instead of as a single number. Because autotuning firmware can hold a miner at any power level in that range, every power level has its own J/TH and the single operating point of stock firmware opens into a curve. \emph{The curves in this book are fleet averages of aggregated operating data from miners running Braiins~OS, second quarter of 2026, not laboratory curves. An individual machine sits somewhere around them, not on them.}
\item[efficiency (J/TH)] \hypertarget{glossary:efficiency}{}A mining machine's power draw divided by its hashrate, in joules per terahash. Lower is better. It is the single most important specification of a mining machine, because it sets the electricity price at which that machine stops being profitable. \emph{This is the one unit in the book where the good direction is down: 23.8~J/TH is more efficient than 29.5. Anchors: the Antminer S9 needed roughly 100~J/TH a decade ago, the S21~XP of 2024 needs 13.5, the latest hydro machines below 10, and this book puts the network's fleet average between 18 and 26~J/TH at the machines as of August 2026, with my own estimate at 21.5~J/TH; the listed fleets that publish their efficiency averaged 17~J/TH, weighted by hashrate.}
\item[Emergency Pricing Program (EPP)] \hypertarget{glossary:emergency-pricing-program}{}The circuit breaker added under the system-wide offer cap in December 2023: once real-time prices have been at \$5,000 per megawatt-hour for 12 hours within any 24, the cap on energy and ancillary-service offers falls to \$2,000 for at least a day and until a day after any emergency ends. \emph{Under it a week like Uri's could not recur at its price, so the values the book measures in that week were priced under 2021's rules. See system-wide offer cap, Winter Storm Uri.}
\item[Emergency Response Service (ERS)] \hypertarget{glossary:emergency-response-service}{}ERCOT's contracted reserve of last resort, deployed when operating reserves fall below 3,000~MW or an emergency is declared, that is, after the ancillary services of Chapter~7 have already been used. A load contracts months ahead to act on ERCOT's call and is paid for the promise whether or not the call ever comes. \emph{Budget capped by rule at \$75 million a year plus \$25 million for emergency renewals; procured four times a year in ten-minute and thirty-minute varieties with an offer cap of \$80 per megawatt per hour; deployment bounded at 24 cumulative hours in a winter contract period and 12 in the others; eleven deployments in eighteen years. It cleared at about \$4.36 per megawatt per hour on average in 2025 and \$18.69 in its summer-afternoon block, which collides directly with the peak layer (see ERS time period).}
\item[energy-only market] \hypertarget{glossary:energy-only-market}{}ERCOT's market design, in which a generator earns money by selling energy and ancillary services and nothing else. If those revenues do not justify building a plant, the plant does not get built. \emph{The design allows prices during scarcity to rise very high, which is exactly what makes the market pay for a load that can get out of the way.}
\item[energy storage resource (ESR)] \hypertarget{glossary:energy-storage-resource}{}ERCOT's name for a grid battery, the one asset class whose entire business is deleting price extremes: it buys the trough and sells the peak. \emph{Every megawatt of it lifts the floor and shaves the spike the miner's option was written on. A battery is duration-limited, about 1.8 hours average storage in 2025, and gives back about 85 percent of what it takes in, so it buys about 1.2 megawatt-hours for every one it sells; a pausing miner has no state of charge, owes no energy back, and can stay off for an hour or a week.}
\item[ERCOT] \hypertarget{glossary:ercot}{}The Electric Reliability Council of Texas, operator of the Texas Interconnection and the market this book is about. It serves about 90 percent of Texas electric load and more than 27 million customers, and because its wires stay almost entirely inside Texas its wholesale market is exempt from FERC rate regulation. \emph{Federal reliability standards still apply through NERC; only wholesale rate regulation is exempt. More Bitcoin mining load is connected here than anywhere else in the world.}
\item[ERCOT Contingency Reserve Service (ECRS)] \hypertarget{glossary:ercot-contingency-reserve-service}{}The ten-minute reserve: capacity ERCOT can call to restore the crash barrier after it has been used, sustained for at least an hour (two hours until RTC+B cut it in December 2025). Launched June 10, 2023, the first new daily procured service in more than twenty years. \emph{Its first year is the book's cautionary case: the market monitor found that the way ECRS was procured and withheld created artificial shortage pricing and added more than \$12 billion of real-time market costs, for a computed reliability benefit of about \$12 million against more than \$600 million of procurement cost. It is the product at the center of the reserves-or-strike trade-off in Chapter~9.}
\item[ERS time period] \hypertarget{glossary:ers-time-period}{}The standing block of hours a load sells into ERS and must stay available for the whole contract term. There are eight: five weekday daytime and evening blocks, two for weekends and holidays, and one for everything left over. \emph{The fourth block is weekdays from 4 to 7 in the afternoon, and 29 of the 32 official 4CP intervals fall inside it, which is why a mine that intends to dodge the peak should never sell that block.}
\item[failure to provide] \hypertarget{glossary:failure-to-provide}{}Under the pre-December 2025 rules, what happened when a committed load stopped consuming: the reserve it had sold vanished, ERCOT bought replacement capacity, and the seller was billed at replacement cost. \emph{Pausing during a \$3,000 spike, the move Chapters~5 and 6 spend forty pages perfecting, was a failure to provide. The charges and the replacement-procurement market were deleted from the rulebook at RTC+B go-live. The two-failure rule still stands: two failures in a rolling year to answer a deployment disqualify a relay-only load, with a six-month wait to reapply. It counts missed deployments, not pauses as such, and it never applied to a CLR.}
\item[fast frequency response (FFR)] \hypertarget{glossary:fast-frequency-response}{}The fastest sub-form of Responsive Reserve, supplied almost entirely by batteries but open to any resource fast enough, a load included: full output within a quarter second of a frequency trigger, sustained at least fifteen minutes. \emph{FFR and the loads' under-frequency relay share one cap: together they may cover no more than 60 percent of the RRS requirement.}
\item[fence (the fence)] \hypertarget{glossary:fence}{}The book's shorthand for a mine behind the meter and for the deal that comes with it, as in Chapter~9's ``The Other Offer: Behind the Fence''. The fenced site pays one negotiated price per megawatt-hour and nothing else: no losses, no ancillary or uplift allocation, no distribution line, no transmission charge. \emph{It also sells no reserves, dodges no peaks and arbitrages no index. The fence is an annuity; the stack is an option. See Private Use Network, indifference price, hybrid.}
\item[FERC (Federal Energy Regulatory Commission)] \hypertarget{glossary:ferc}{}The U.S. regulator of interstate electricity transmission and wholesale sales. See NERC for the organization that develops and enforces bulk power system reliability standards.
\item[firm load] \hypertarget{glossary:firm-load}{}A load contracted to run whatever the price does. Today's AI campus is built and contracted as one, because nobody pauses a training run for a price. \emph{The opposite of the flexible layer beside it: in Winter Storm Fern ERCOT credited about 4,100~MW of curtailment to miners against about 170~MW to data centers. Not the same use of ``firm'' as firm supply, which is a duty on the generator.}
\item[firm supply] \hypertarget{glossary:firm-supply}{}A behind-the-meter deal in which the generator must keep the power flowing, so the mine runs whenever it wants. \emph{The firm case is a ceiling for the fence: an always-available plant is the most favorable assumption the ledger makes for it. Contrast interruptible.}
\item[flexibility] \hypertarget{glossary:flexibility}{}The product a mine sells to the grid: the freedom to change how much power it draws inside a quarter-hour, on a price signal, a demand signal or a dispatch instruction, at a cost the machines were built to bear. \emph{In the ledger the word also names a quantity: the planned site's margin over the never-pausing site's, worth 9.0, 8.0 and 15.8 percent of the never-pausing margin in 2022, 2024 and 2025 and 23.0 and 29.5 percent in 2021 and 2023, the storm year and ECRS's first year. The best annual choice in hindsight would have made the storm year 52.6 percent; see hindsight ceiling. Every such margin is after power and grid charges; see never-pausing site.}
\item[the flexibility stack (the stack)] \hypertarget{glossary:the-flexibility-stack}{}The layers a site runs at once to earn from flexibility: the strike, the power target and the summer window, with reserves beside them. They are not additive, because a megawatt does one thing in one quarter-hour. \emph{Chapter~9's own correction to its title: the strategy is not a stack but an order of switching. Inside the year it runs with no look-ahead at all; its two annual switches, reserves and the summer window, are the only choices the ledger could have made with hindsight, and its headline sets them in advance (see planned site, hindsight ceiling). On my reading, under RTC+B and, if it is adopted as proposed, 12CP, the stack collapses toward one layer, the moving strike with the power target beside it, and falling hashprice makes that layer thicker.}
\item[flexible load] \hypertarget{glossary:flexible-load}{}A load that can rapidly scale down or shut off its power usage when the grid needs it to. Bitcoin mining is the book's central example. \emph{ERCOT's Winter Storm Fern accounting put roughly 4,100~MW of the January 2026 reduction on cryptocurrency miners, second only to the oil and gas sector and more than twenty times what conventional data centers contributed.}
\item[frequency deadband] \hypertarget{glossary:frequency-deadband}{}The band of frequency error inside which no response is demanded. A Controllable Load Resource's may be no wider than plus or minus 0.036~Hz against a 60.000-hertz target, so outside that band it must answer like a power plant's governor (under a change ERCOT schedules for early 2027, only if it is qualified for Regulation or Responsive Reserve). \emph{Not the price deadband of Chapter~5, which is a band of dollars, not hertz. See droop, Controllable Load Resource, price deadband.}
\item[generation resource] \hypertarget{glossary:generation-resource}{}The full ERCOT resource class, which may offer energy, be committed and dispatched, and sell ancillary services. In exchange it takes on continuous telemetry, a standing offer whenever it is online, and the duty to run when ERCOT commits it for reliability. \emph{The counterpart to a Settlement Only Generator; the choice between the two is the gate that matters for a miner who owns a plant.}
\item[halving] \hypertarget{glossary:halving}{}The scheduled event, about every four years, that cuts the block subsidy in half and so reduces the issuance rate of new bitcoins. \emph{Four so far: 2012, 2016, 2020, 2024. The stretch between two of them is a halving epoch, the book's only sense of the word epoch.}
\item[halving epoch] \hypertarget{glossary:halving-epoch}{}The stretch between two halvings, about four years, across which the block subsidy is constant. \emph{In this book epoch always means a halving epoch. The protocol's other recurring interval, the two-week window between difficulty adjustments, is called a difficulty period here and never an epoch. See halving, difficulty adjustment.}
\item[hashprice] \hypertarget{glossary:hashprice}{}The dollars a unit of hashrate earns per day: network revenue per day, subsidy plus fees, times the bitcoin price, divided by the network hashrate. The book quotes it in dollars per terahash-day (\$/TH/day), a terahash-day being one terahash per second of hashrate sustained for a day. \emph{Three principal drivers shape it: the bitcoin price, the network hashrate and the scheduled halving. The book's anchor year is 2025, which averaged about \$0.051 per terahash-day. Watch the units: market convention also uses \$/PH/day, a thousand times larger.}
\item[hashrate] \hypertarget{glossary:hashrate}{}The computational power a miner or the whole network brings to the puzzle, measured in hashes per second: terahashes per second (TH/s) for a machine, exahashes per second (EH/s) for the network. \emph{Network hashrate is about 930~EH/s as of August 2026, down from a peak of roughly 1,100~EH/s in late 2025.}
\item[hindsight ceiling] \hypertarget{glossary:hindsight-ceiling}{}The best annual choice in hindsight: the site runs on whichever reserve allocation, product and summer-window call scored best on the year's own outcomes. \emph{The oracle's discipline applied one level up, to the annual switches; inside the year it still runs the lagged loop. Printed beside the planned site and never called a strategy, because it needs the year's prices before the year starts; in 2021 it needs to know in January that Uri is coming. Its flexibility was 52.6 percent of the never-pausing margin in 2021 and \$233 million over the five years, against the planned site's 23.0 percent and \$123 million.}
\item[hosting agreement] \hypertarget{glossary:hosting-agreement}{}The contract under which a site's owner, the host, runs a tenant's machines on its power. It decides who holds the right to curtail. \emph{In the public Rockdale agreements of Chapter~6 the host holds it: one has the host sell the power instead once the ERCOT market's profitability passes a redacted share of the contract price, the other lets it enroll the site in any demand response program at its sole discretion, and each passes part of that value back to the tenant. The book's site is an owner-operator's, which buys its own power and holds that right itself.}
\item[hub] \hypertarget{glossary:hub}{}A trading settlement point whose price averages the LMPs of a fixed list of major 345 kV buses. Traders hedge at hubs (North, South, West, Houston and a few others), and hub prices anchor most financial contracts. \emph{A hub is a financial construct, not a place a load can be billed at; a consumer still pays its load zone, and the gap between the two is the basis.}
\item[hybrid] \hypertarget{glossary:hybrid}{}A site that sits behind a generator and in front of the grid at once, keeping both the fence's discount and its own switch. \emph{The hybrid buys the annuity and keeps the option. The book describes it but does not model it, because the companion repository holds no generator-availability data. Hut 8's Vega runs behind a wind farm and in front of the grid at once. Not a generator co-located with a battery, which the trade press also calls a hybrid.}
\item[hydro] \hypertarget{glossary:hydro}{}Shorthand for a liquid-cooled mining machine. Hydro models sit at the efficiency frontier, with the latest pushing below 10~J/TH. \emph{Not hydroelectric power. In mining, hydro describes how the machine is cooled, not where its electricity came from; Chapter~1 mentions hydroelectric dams in Canada, which is the other sense.}
\item[ICCP] \hypertarget{glossary:iccp}{}The grid world's standard control-room protocol: the language a QSE's control system uses to speak to ERCOT across ERCOT's own private wide-area network, on routers ERCOT supplies and manages at the QSE's premises. \emph{ERCOT's side of the private-network hookup is capped at \$7,000 to install and \$450 a month to run.}
\item[Independent System Operator (ISO)] \hypertarget{glossary:independent-system-operator}{}A neutral operator that runs the transmission system and the wholesale market on top of it. ISOs and RTOs own no generation and serve no retail customers: they dispatch generators, operate the transmission system and run the auctions where electricity is bought and sold. \emph{Roughly two-thirds of U.S. electricity consumers live in an ISO or RTO region; seven such markets serve them, and ERCOT is one. The book uses ISO and RTO for the same kind of institution.}
\item[index contract] \hypertarget{glossary:index-contract}{}A power contract under which the buyer pays the market's settlement price for every interval it consumes. The buyer carries all of the price risk. \emph{Texas banned wholesale-indexed products for residential and small commercial customers after Uri and left them fully legal for large loads. If you index, index real time (see real-time index).}
\item[indifference price] \hypertarget{glossary:indifference-price}{}The negotiated behind-the-meter price at which the fenced site earns exactly what the grid site earns. Below it the fence wins, above it the grid does. \emph{Quoted against three grid sites and the answers differ: against the planned site it was \$55, \$63, \$43, \$41 and \$44 from 2021 to 2025; against the hindsight ceiling it had no positive answer in 2021; against the never-pausing site it was \$192, \$84, \$76, \$50 and \$56.}
\item[inertia] \hypertarget{glossary:inertia}{}The energy stored in the spinning mass of every synchronous generator locked in step with the grid, which the grid draws on in the first seconds after a plant trips. It slows the fall of frequency; it cannot stop it. ERCOT's typically runs between 200 and 400 gigawatt-seconds, against a floor of 100. \emph{On ERCOT's accounting, wind, solar and batteries, which reach the grid through inverters, add none. During Uri's frequency emergency it stayed between 254 and 349 gigawatt-seconds: that night's danger was a shortage of reserves, not of inertia.}
\item[interconnection] \hypertarget{glossary:interconnection}{}A vast region whose generators all run synchronously at the same electrical frequency, 60~Hz in the United States, so power can flow across regional lines as needed. North America has the Eastern, Western and Texas interconnections. \emph{Texas is its own interconnection, which is why ERCOT can neither lean on nor be dragged down by its neighbors the way the Eastern Interconnection spread the 2003 cascade.}
\item[interference] \hypertarget{glossary:interference}{}How far the stack falls short of adding its layers separately: the joint price-only result less the naive sum of strike, power target and the whole transmission prize, negative in every year. \emph{The overstatement was 15 percent in 2025 and 34 percent in 2024, when two of the four peaks also escaped the summer window; in dollars it was largest in the storm year, because the strike was highest then. It exists because a megawatt does one thing in one quarter-hour.}
\item[interruptible] \hypertarget{glossary:interruptible}{}A behind-the-meter deal in which the generator keeps the right to take the power back and sell it to the grid, paying the mine nothing for the interruption. \emph{An interruptible fence matches the planned site at no price in any of the five years. Firmness matters alongside price in a behind-the-meter deal. Distinct from interruptible load, which is a property of the load, not a contract term.}
\item[interruptible load] \hypertarget{glossary:interruptible-load}{}A load that can be stopped for hours without ruining a product or inconveniencing anyone. Stopping a miner only delays hashing, which is why mining is interruptible by design. \emph{Distinct from interruptible, which is a contract term of Chapter~9, not a property of the load.}
\item[islanded] \hypertarget{glossary:islanded}{}Running with no grid tie at all. An islanded mine avoids every non-energy charge and loses the whole flexibility stack with it. \emph{It has no load zone price to arbitrage, cannot sell reserves, has no transmission peak to dodge and cannot enroll in the emergency-response programs. Distinct from a netted site, which keeps a grid tie and with it part of the flexibility stack.}
\item[Large Computational Load] \hypertarget{glossary:large-computational-load}{}A Large Load where half or more of the demand is computational load, such as data centers and cryptocurrency mining facilities. \emph{This is the term the commission's transmission reform asks about by name.}
\item[Large Load] \hypertarget{glossary:large-load}{}ERCOT's name for a load of 75~MW or more at a single site. It must complete ERCOT's large-load interconnection process, with system studies and a separate approval to energize, before it may draw power. \emph{By the end of 2025 ERCOT's large-load program, which takes loads of every industry, had about 5,800~MW energized, each site counted at its highest draw. By March 2026 over 9,000~MW, including those, held approval to energize; mining demand alone peaked near 4,600~MW in 2025 (Chapter~7). The 75~MW line is also where the proposed large-load demand management service would apply and, for a load that connects or expands under the commission's new interconnection rule, the proposed minimum billing demand.}
\item[large-load demand management service] \hypertarget{glossary:large-load-demand-management-service}{}The emergency product Senate Bill 6 adds alongside ERS, for sites of 75~MW or more called with 24 hours' notice. It bars any load that already curtails on price or sells another reliability or ancillary service, so it is written for the loads that do not already respond. \emph{The rule is still a proposal, which the commission's staff plans to bring for adoption in January 2027, so the service is not yet running; the proposal caps offers at \$5,000 per megawatt-hour and extends the exclusions to ERS and the utility programs.}
\item[load] \hypertarget{glossary:load}{}The electricity being consumed by end users at any given moment. The book uses load and demand for the same thing. \emph{A mining site is a load, never a generator, even when switching it off has the same balancing effect as adding supply.}
\item[Load Frequency Control (LFC)] \hypertarget{glossary:load-frequency-control}{}ERCOT's name for automatic generation control (AGC): the control-room system that, in the minutes after primary frequency response has stopped a fall, moves the generators and controllable loads selling Regulation until frequency is back at 60~Hz, sending a correction signal every four seconds. \emph{It brings frequency home; primary response only stops it falling.}
\item[load resource] \hypertarget{glossary:load-resource}{}A load registered with ERCOT to sell ancillary services. In its basic form it is wired to trip offline within half a second on an under-frequency relay, or curtailed on ERCOT's instruction, and paid for it. The ability to disconnect quickly is itself a sellable service. \emph{The Controllable Load Resource is the other form of the same registration: instead of tripping or dropping on a call, it follows a base point.}
\item[load shedding (firm load shed)] \hypertarget{glossary:load-shedding}{}The grid operator deliberately cutting power to blocks of customers, in rolling blackouts, to force supply and demand back into balance and save the rest of the system. ERCOT's and Senate Bill 6's name for it is firm load shed. \emph{Reserved for involuntary outages ordered by the operator (the automatic form, set off by relays with no order, is under-frequency load shedding): never a substitute word for curtailment, and never for a miner's pause. Customers who are shed do not choose and are not paid. During Uri, 1,000~MW was ordered at 1:20 a.m., 5,000~MW by 1:50, 8,500~MW by 1:55 and a maximum of 20,000~MW later that day. Under SB 6 any transmission-connected load interconnected after 2025, other than critical-load industrial customers and critical natural gas facilities, must be curtailable during it, under a protocol its wires utility puts in place before the load connects; the duty is written into the statute but not yet into a commission rule.}
\item[load zone] \hypertarget{glossary:load-zone}{}The settlement point a consumer pays: an average over every bus in the zone, weighted by how much load each bus is carrying. ERCOT defines eight, the four large competitive zones plus four belonging to utilities that did not open to retail competition (Austin Energy, CPS Energy, LCRA and Rayburn). \emph{Zones are sets of electrical buses, not geographic borders. Until nodal settlement for Controllable Load Resources switches on, every consumer's energy is billed at the zonal price.}
\item[Locational Marginal Price (LMP)] \hypertarget{glossary:locational-marginal-price}{}The cost of delivering one more megawatt-hour at exactly one node, computed by SCED roughly every five minutes at every node in the network model, and in the DAM once a day for each hour of the next operating day. In ERCOT it has two components only: the System Lambda plus the Congestion Cost. \emph{Unlike PJM, MISO and the other nodal markets, ERCOT's LMP has no loss component. Losses are recovered by scaling metered consumption up by a loss factor, so a site's location changes its price only through congestion, not through distance from the plants.}
\item[the loop] \hypertarget{glossary:the-loop}{}The book's name for the price-response layer of Chapter~5: it reads each new print, recomputes the profit-optimal power target, and applies the hold, the price deadband and the backstop. It owns both decisions a fleet makes on price, but never in one rule: modulation sets the power target, the backstop owns pause and resume, and the two read the same feed without sharing a decision. \emph{An abstraction on purpose, so the design can be argued on its own terms rather than any one vendor's. Braiins Manager's Price Adapt implements the loop (Chapter~5). Below the loop sits the autotuning firmware, which holds whatever target the loop asks for unless temperature safety forces it lower. See modulation, backstop, power target, autotuning firmware.}
\item[LZ\_WEST] \hypertarget{glossary:lz-west}{}ERCOT's West load zone, the one that appears throughout this book. West Texas has the state's best wind and solar resources, limited transmission out and much of the state's Bitcoin mining, which gives it the most volatile prices in ERCOT. \emph{Its 2025 average 15-minute real-time settlement point price, \$42.18 per megawatt-hour, is the book's price anchor.}
\item[marginal machine] \hypertarget{glossary:marginal-machine}{}The least efficient machine still running on the network at a given hashprice. Because the fleet lives on the breakeven line by construction, the marginal machine mines for nearly nothing, and its only remaining edge is the price it pays for power. \emph{Defined by the market, not by a model: it is whichever machine sits at the breakeven line today. This is why only the cost side decides who survives.}
\item[market monitor] \hypertarget{glossary:market-monitor}{}Potomac Economics, the independent firm that watches ERCOT's wholesale market for market power and design flaws and reports what it finds to the commission. The commission selects and supervises it, and ERCOT funds it out of the system administration fee that pays for ERCOT itself. It holds no enforcement power: statute reserves all of that to the commission. \emph{Statute and rule call it the independent market monitor, or IMM. Its annual State of the Market report is a standing source here, behind the ECRS verdict of Chapter~7 and the 4CP criticism of Chapter~8; the bibliography lists it under Potomac Economics.}
\item[mean reversion] \hypertarget{glossary:mean-reversion}{}The tendency of five-minute prices to fall back after a large upward leap rather than keep climbing; after a large fall they usually keep falling, as a spike unwinds over several prints. The level, not the trend, carries the information. \emph{At LZ\_WEST, from 2024 to mid-2026, the correlation between one five-minute price change and the next is about -0.29, and the spikes carry it: each of the five largest leaps, all of more than \$4,000, fell straight back at the very next print, while ordinary moves of a few dollars carry almost no signal either way, and the first half of 2026, when no print leapt that far, came out near zero. This is the structural reason there is no forecast in the loop: extrapolating a trend points the decision the wrong way, and it is most wrong during spikes, when the stakes are highest.}
\item[minimum billing demand] \hypertarget{glossary:minimum-billing-demand}{}Under the proposed rule, the demand billed whether or not the site consumes, for a load the commission's new large-load interconnection rule covers: a new site of 75~MW or more, or the added load of an expansion that first takes a site to 75~MW or adds 75~MW or more to one already past it. It is the greater of the peak demand written into its service contract, its highest demand in the past year, or its 12CP demand, billed for at least twenty years. \emph{As proposed, for such a load this ends peak avoidance altogether: transmission becomes a fixed cost. A site already energized when that rule takes effect on October 8, 2026, and not expanding, is outside it as proposed. The chairman of the commission called it a minimum demand ratchet.}
\item[mining pool] \hypertarget{glossary:mining-pool}{}A service that combines miners' computing power and shares the proceeds.
\item[modulation] \hypertarget{glossary:modulation}{}The gradual repricing of a running machine: moving its power target up and down the efficiency curve as the electricity price changes. \emph{Modulation never pauses or resumes anything; the backstop owns that. The two read the same price feed and never share a decision, which keeps noise out of the pause and resume path.}
\item[modulation floor] \hypertarget{glossary:modulation-floor}{}The lowest power target modulation will set: Braiins~OS's default minimum, or the machine's best-efficiency point where that is higher, since below the best point slowing further surrenders more hashrate than the electricity it saves. In this book, 1,200 W for the S19j~Pro and 2,774 W for the S21~XP. \emph{Past the floor there is only one move left, the full pause, and that decision belongs to the backstop. The loop clamps every computed target between this floor and the power ceiling.}
\item[negative price] \hypertarget{glossary:negative-price}{}An hour in which generators pay consumers to take power, because subsidized renewables can profitably produce below zero and shutting down costs them more. In West Texas it happens when renewable output exceeds what the wires can carry away. \emph{At a generator's own node the real-time price has a floor of -\$251/MWh. At LZ\_WEST in 2025 prices were negative 5.6 percent of the time.}
\begin{samepage}
\item[NERC (North American Electric Reliability Corporation)] \hypertarget{glossary:nerc}{}The organization that develops and enforces reliability standards for the North American bulk power system. See FERC.\par
\end{samepage}
\item[net load] \hypertarget{glossary:net-load}{}Demand minus wind and solar output. Its highest hour, the net-load peak, is the hour the system is actually short of power. \emph{ERCOT put the net-load peak at hour ending 21 in 2023 and 2024 and hour ending 22 in June and July 2025: the tightest hours of the day have migrated from the traditional four-to-five-p.m. demand peak to the seven-to-nine-p.m. net-load peak, the hour solar leaves. Not the 4CP interval, which counts gross demand with solar still shining and has stayed at 5 p.m.; ERCOT told its board in 2026 that 4CP does not correspond to the system's highest net-load peaks.}
\item[netting] \hypertarget{glossary:netting}{}Settling one number at the interconnection point in every quarter-hour, generation minus load. Net export settles as generation, net consumption settles as load and carries the load-side charges, and consumption that never crosses the meter is invisible to settlement. \emph{ERCOT forbids netting by default and allows it in five enumerated site configurations. Since Senate Bill 6, a load of 75~MW or more that seeks a new or expanded connection and nets against a plant already registered as stand-alone on September 1, 2025 needs the commission's approval first, on a 120-day ERCOT study and a decision 60 days after that, approved by default if the deadline passes.}
\item[never-pausing site] \hypertarget{glossary:never-pausing-site}{}The comparator site: the same fleet running at stock power in every interval, present at all four summer peaks, and paying every grid charge in full. Every gain in the ledger is measured against it. \emph{A baseline, not a strategy, and the denominator of the flexibility percentages. Its figure is a margin after power and grid charges; machines, depreciation, pool and firmware fees, site operation, staff and debt are outside the ledger.}
\item[node] \hypertarget{glossary:node}{}A physical junction where equipment connects to the high-voltage grid, such as a plant's switchyard or a substation where lines meet and transformers step the voltage down. The power network calls it a bus; the market calls it a node, and a large mine's own interconnection point is one. \emph{Bus and node are the same thing under two names; a node is a real place with coordinates, not a zone or a region. Not a Bitcoin network node, which is a computer running the protocol.}
\item[Non-Spinning Reserve (Non-Spin)] \hypertarget{glossary:non-spinning-reserve}{}ERCOT's slow reserve: capacity that can come online within thirty minutes, with four hours of endurance (two hours once the Dispatchable Reliability Reserve Service is running), the service that covers a forecast bust rather than an instant failure. \emph{The only product whose standing revenue was firming rather than falling at the 2026 run rate.}
\item[off-peak nights (7x8)] \hypertarget{glossary:off-peak-nights}{}The eight hours outside the daytime block, midnight to 6 a.m. and 10 p.m. to midnight, every day of the week, weekends and holidays included. \emph{One of the three standard block products, and 2,920 hours in 2025. At LZ\_WEST that year the window averaged \$49 per megawatt-hour and held firm while the daytime window sagged, because it is the one window solar never reaches. The trade also says off-peak for everything outside 5x16, which is 7x8 and 2x16 together; this book always means the nights. See block.}
\item[off share] \hypertarget{glossary:off-share}{}The share of quarter-hours in which a fleet draws nothing under the price-only stack, where the power target underclocks the machines before anything stops. \emph{Not the pause share. The S19j~Pro that paused 14 percent of 2025 under the plain strike rule drew nothing in only 7 percent of quarter-hours under the stack, because the engine underclocks first.}
\item[on-peak (5x16)] \hypertarget{glossary:on-peak}{}The sixteen hours from 6 a.m. to 10 p.m. Central time, Monday through Friday, excluding six market holidays. The market labels each hour by when it ends, so the window runs HE 7 through HE 22. \emph{One of the three standard block products, and 4,080 hours in 2025. At LZ\_WEST that year the window averaged \$41 per megawatt-hour, cheaper than the \$49 nights, because it now contains the solar belly. See block, 2x16, off-peak nights (7x8).}
\item[operating leverage] \hypertarget{glossary:operating-leverage}{}The sensitivity of a machine's margin to a move in hashprice, which grows as the share of revenue spent on electricity grows. \emph{On the eve of the August 2026 rally, at the \$42.18 price anchor, the S19j~Pro was spending 96 cents of every revenue dollar on electricity and the S21~XP about 44, so a ten percent move in hashprice barely registers at the efficient machine and doubles or erases the inefficient one's entire margin. Efficiency raises profits and reduces exposure to hashprice.}
\item[Operating Reserve Demand Curve (ORDC)] \hypertarget{glossary:operating-reserve-demand-curve}{}ERCOT's administered scarcity pricing mechanism from 2014 until the 2025 redesign: a real-time adder to the energy price that grows as available reserves shrink, reflecting the rising probability of losing load. When reserves were plentiful the adder was near zero. \emph{RTC+B retired the adder on December 5, 2025: shortage value now sits inside the energy price rather than being added to it afterwards, and only the curve's shape survives, in the ancillary-service demand curves cut from it.}
\item[operating reserves] \hypertarget{glossary:operating-reserves}{}The grid's safety margin, capacity held back and ready rather than producing. \emph{Valuable even in hours when no shortage materializes, which is why ERCOT pays for it administratively. Chapter~2 introduces the concept; Chapter~3 prices it.}
\item[optimal site] \hypertarget{glossary:optimal-site}{}Not a term of this ledger. Chapter~9 prints the planned site as its headline and the hindsight ceiling beside it; neither is optimal for a site that cannot see the year ahead. \emph{Chapter~10's sensitivity rows, which rerun 2025 alone, use the phrase for the price-only stack with the summer window played where it pays, which in 2025 is also the planned site. See planned site, hindsight ceiling.}
\item[oracle (perfect-foresight twin)] \hypertarget{glossary:oracle}{}The benchmark the loop is graded against: a hypothetical operator with perfect foresight and zero friction. In every interval it sets the machine to the profit-optimal power for that interval's own settlement price, which no one can know until the interval has ended, and it pauses for free whenever even that optimum would lose money. No hold, no price deadband, no minimum runtimes, no resume ramp, no publication delay. \emph{A ceiling on what responding to the energy price can earn, not a strategy; nothing real sees the future. Not a price oracle in the crypto sense: nothing here feeds data to a chain. Run alongside the live loop as a perfect-foresight twin so every backtest can state the price of not seeing the future. In the LZ\_WEST backtests the full loop captures 95 to 98 percent of the oracle's gain over stock power on the newer fleet and 93 to 96 on the older one. Chapter~9's simpler twin only pauses, at stock power: against it the S19j~Pro's plain strike, decided one settlement interval late, took 89 percent of the twin's gain over never pausing on the real-time index in 2025.}
\item[pass-through] \hypertarget{glossary:pass-through}{}The third layer of an index bill: delivery charges, 4CP transmission and ancillary service costs billed by others and handed to the miner unchanged. \emph{Where first-time buyers get surprised, because none of it appears in the quoted price. The Rhodium contract makes the customer sign a specific acknowledgment that in a bad month ancillary charges alone can multiply the bill.}
\item[pause] \hypertarget{glossary:pause}{}The machine's own act of stopping hashing: past the modulation floor there is only one move left, the full pause, and the loop hands that decision to a separate rule, the backstop. \emph{A pause is curtailment of the self-initiated kind: the narrower word, used where it matters that the decision was the site's own rather than an instruction it was sent. The book never substitutes load shedding, powering down or demand reduction for either. See curtailment.}
\item[pause share] \hypertarget{glossary:pause-share}{}The share of a period's quarter-hours in which the plain strike rule stopped a fleet. \emph{For the S19j~Pro it rose from 2 percent of 2021 to 14 percent of 2025, about 1,200 hours, each year scored at its own strike, which fell from about \$440 to \$72 as hashprice fell: the rise is hashprice, not ERCOT getting spikier. Distinct from off share, which is the joint engine's total off time and includes the hours it would otherwise have underclocked.}
\item[persistence] \hypertarget{glossary:persistence}{}Using the last published print as the decision price, with no forecast, no smoothing and no trend extrapolation. It is a measured choice, not a shortcut. \emph{Sixty-seven alternatives to it were tested, spanning slope extrapolation over 10-to-60-minute windows, exponential smoothing, rolling quantiles and constant risk biases. Out of sample none beat it by anything worth building: on the S19j~Pro fleet nothing beat it at all, on the S21~XP the best edged ahead by 0.2 percent of the plain hold's gain, and across the eight load zones only the two gentlest slope settings came out ahead, by at most 0.12 percent. Leaning the decision price against the recent move does little better.}
\item[planned site] \hypertarget{glossary:planned-site}{}Chapter~9's headline: the site with its two annual switches set as an operator could have set them in advance. In January it commits to the reserve allocation and product that won the previous year's ledger (no reserves in 2021, the first year of the data); on June 1 it calls the summer window on yesterday's hashprice and last summer's window price. \emph{Flexibility in the ledger is the planned site over the never-pausing one. Inside the year every decision is lagged, never foreseen. Its flexibility, \$123 million over the five years, sits inside the \$112 to \$127 million of the rules an operator could have set in advance; never selling reserves earned the most, since every rule set in advance that said sell lost to it, but the plan stays the headline because its rule was chosen for what it knew, not for what it earned. Its first January has no previous year to follow, so it takes the default that needs no history, no reserves; had 2021 gone into RRS, as its own year rewarded, its flexibility over the five years would come to \$227 million, by luck rather than by plan. Distinct from the hindsight ceiling, which picks the same switches with the year's own outcomes.}
\item[postage stamp rate] \hypertarget{glossary:postage-stamp-rate}{}The one ERCOT-wide rate, in dollars per kilowatt of 4CP demand, that every transmission owner's individual rate adds up to: the same price for a kilowatt of peak demand anywhere in the state. \emph{For 2025 it was \$68.55 per kilowatt-year on a 4CP of 81~GW, recovering a \$5.4 billion cost of service; because it is the sum of each utility's own rate, it sits about 2 percent above that cost divided by the 4CP. This is the wholesale number; what a site actually pays is the retail rider its own wires utility charges, which for a transmission-voltage site landed below it on the riders of September 2025.}
\item[power ceiling] \hypertarget{glossary:power-ceiling}{}The highest power target the loop will set: the top of its power range, 3,420 W for the S19j~Pro and 3,990 W for the S21~XP in this book's backtests, below the highest power targets the fleet dataset records. Cheap hours push the profit optimum up against it. \emph{The counterpart of the modulation floor; the loop clamps every computed target between the two. It sits above stock power, because overclocking buys extra hashrate at ever-worse efficiency and only cheap electricity pays for that trade.}
\item[power generation company] \hypertarget{glossary:power-generation-company}{}The commission registration class a generator takes if it means to sell its output, whatever its size. Registration is required before the first day of generation. \emph{The first of the four gates a mine must pass to sell a megawatt-hour into ERCOT. There is no megawatt floor in this class.}
\item[power purchase agreement (PPA)] \hypertarget{glossary:power-purchase-agreement}{}A bilateral contract that fixes the volume, price and term of the power bought. The supplier carries the market risk and charges for taking it. \emph{A PPA without a sell-back clause kills the pause option: the power is paid for whether you use it or not, so its marginal cost to you is zero and no market price can make pausing profitable.}
\item[power target] \hypertarget{glossary:power-target}{}The power level a machine is told to hold, recomputed every five minutes from the newest print and clamped between the modulation floor and the power ceiling. The autotuner then reaches it, either by retuning every chip's voltage and frequency or by applying a profile it already holds for exactly that target, since a saved profile serves no other. \emph{This is the layer that overclocks or underclocks; it is not the pause layer. See the loop, autotuning firmware.}
\item[power-weighted strike] \hypertarget{glossary:power-weighted-strike}{}The site-wide breakeven price: a strike computed as any strike is (see strike price), but on the site's power-weighted efficiency, total watts over total terahash, which for this fleet is 17.95~J/TH. It is the price the summer window rule sets against last summer's window price. \emph{Do not use it to control pauses: a loop using it would pause the S21~XP when only the S19j~Pro should stop and run the S19j~Pro when it should already be dark. Each model pauses on its own strike.}
\item[price arbitrage] \hypertarget{glossary:price-arbitrage}{}Acting on the gap between what a megawatt-hour costs you and what it is worth to you: consume when it is worth more than it costs, and pause, or sell the power back, when it costs more than it is worth. \emph{The stopping here is the site's own decision, so this book calls it a pause.}
\item[price deadband] \hypertarget{glossary:price-deadband}{}The one extra condition that breaks the loop's 30-minute hold: re-apply the power target immediately whenever the newly computed ideal target diverges from the applied one by more than 10 percent of the machine's default power. \emph{In a quiet market it never fires; on a regime shift the fleet reprices within five minutes instead of thirty. Out of sample it adds 2.6 percent to what the plain hold earns over stock power on the S19j~Pro fleet, at about 26 target moves a day against the hold's 13, and 0.5 percent on the S21~XP at under three; widening it to 25 percent keeps about four-fifths of that on the S19j~Pro while cutting the extra moves roughly in half, and under a quarter on the S21~XP. It beat the plain hold in all sixteen zone-and-fleet combinations, in 2024 and out of sample, at widths of 10 to 25 percent, and on the S19j~Pro it wins back 97 percent of what the plain hold gives up against retargeting on every print. A modulation mechanism only: adding it changed backstop pause counts by exactly zero in every zone. Not the frequency deadband, which is a different quantity under the same word. See frequency deadband, modulation.}
\item[price duration curve] \hypertarget{glossary:price-duration-curve}{}A sort of every interval of the year from most expensive to cheapest, so that the far left shows the crisis hours and the far right the giveaway hours. \emph{LZ\_WEST 2025: about 35,000 15-minute intervals, mean near \$42/MWh, negative 5.6 percent of the time, maximum near \$3,700.}
\item[primary frequency response (PFR)] \hypertarget{glossary:primary-frequency-response}{}The first layer that stops a falling frequency, and the generator's own contribution to Responsive Reserve: turbine governors automatically feed in more steam or gas as frequency falls, without waiting for any instruction, typically starting about three and a half seconds after a trip. It stops the fall and steadies frequency a little above its low point, but does not bring it back to 60~Hz; Load Frequency Control does that. \emph{A Controllable Load Resource must pass a frequency-response self-test to show it can do the equivalent with kilowatts; under a change ERCOT schedules for early 2027, the frequency duty applies only to one qualified for Regulation or Responsive Reserve.}
\item[print] \hypertarget{glossary:print}{}One published real-time price. ERCOT's own documents call it an LMP; the settlement point price is the 15-minute price built from the LMPs. The book borrows the trading floor's word for each published number, and the loop decides on the newest print. \emph{Every print is in force before anyone outside ERCOT can read it. ERCOT's posting lands a median 22 seconds after the five-minute boundary, and the price is readable through its public API a median 2.9 minutes after it, so the backtests model the loop as seeing each print one full interval late, which covered 95 percent of the prints polled.}
\item[Private Use Network (PUN)] \hypertarget{glossary:private-use-network}{}In ERCOT's own words, an electric network connected to the ERCOT Transmission Grid that contains load not directly metered by ERCOT, typically netted with internal generation. In practice: a mine built beside a power plant and metered together with it where the pair meets the grid, so that only the net flow is ever settled. \emph{The trade calls it behind the meter, and this book calls such a site the fence. Chapter~6 prices what such a network is worth to a miner, and what it costs.}
\item[Proof-of-Work (PoW)] \hypertarget{glossary:proof-of-work}{}The consensus protocol that secures the Bitcoin ledger. Specialized computers compete to solve a cryptographic puzzle by repeatedly hashing block data, guessing a random number until the resulting hash falls below a target threshold, and the first to find a valid solution broadcasts it to the network as proof that the work was done. \emph{The work is the energy spent guessing, not any useful computation. This is what makes the security of the network a direct function of electricity consumed.}
\item[protective relay] \hypertarget{glossary:protective-relay}{}The automatic protection that lets a generator disconnect itself to avoid damage when frequency leaves the safe band. In ERCOT, if frequency stays at or below 59.4~Hz for more than nine minutes, generators are allowed to trip themselves offline. \emph{During Uri frequency sat below that line for four minutes and 23 seconds, roughly half the time the relays would have waited; four and a half minutes more and roughly 17~GW would have tripped, risking a blackout of the whole ERCOT grid.}
\item[Provisional Controllable Load Resource (PCLR)] \hypertarget{glossary:provisional-controllable-load-resource}{}The second elective route out of ERCOT's 2026 batch study, for a large load that brings no plant. It may draw above the demand the study says the network can serve, up to its full request and with no guarantee that the grid can serve it in real time, but only while it stays registered as a Controllable Load Resource taking ERCOT's dispatch, with its bids capped when a nearby constraint is at risk and no ancillary services to sell. \emph{The rules passed in June 2026, but the dispatch code behind them is a build ERCOT has ranked for 2027.}
\item[Qualified Scheduling Entity (QSE)] \hypertarget{glossary:qualified-scheduling-entity}{}The intermediary a resource entity must designate to face the ERCOT market on its behalf: a firm it hires or, once it qualifies as one, itself. \emph{It submits the offers, receives the dispatch and settles every payment, and designating one is a registration requirement, not an option; ERS offers too may be submitted only by a QSE, on the load's written authorization. A mine that sells ERCOT nothing still reaches the market through one: the QSE of whoever sells it power. What a QSE charges a site is private, and the book prints no figure. See aggregator, resource entity.}
\item[rate of change of frequency (RoCoF)] \hypertarget{glossary:rate-of-change-of-frequency}{}How fast frequency moves in the first moments after a sudden imbalance, in hertz per second, set by the size of the imbalance and the inertia online. \emph{The less inertia, the steeper the fall and the less time governors, batteries and load relays have to act.}
\item[Real-Time Co-optimization plus Batteries (RTC+B)] \hypertarget{glossary:real-time-co-optimization-plus-batteries}{}ERCOT's December 5, 2025 redesign, its largest since nodal pricing. Every five-minute SCED run now procures energy and ancillary services together and reprices both continuously, instead of buying ancillary services once in the Day-Ahead Market and holding them fixed all day; the same engine that sets base points awards every ancillary product, choosing for each megawatt whether it is worth more as energy or as readiness. \emph{For flexible loads this is a regime change: the value of standing ready is marked to market every five minutes rather than locked in the afternoon before. Day-ahead ancillary awards are now strictly financial positions settled against real-time awards every interval; a load can sell out of its day-ahead position in real time with no failure-to-provide penalty, leaving only imbalance risk. The Day-Ahead Market still does not co-optimize a load's energy purchase against its ancillary offer; that is scheduled for early 2027. The ORDC adder is gone, absorbed into the energy price; only its curve survives, as the shape the new ancillary-service demand curves are cut from.}
\item[real-time index] \hypertarget{glossary:real-time-index}{}A retail contract priced off the real-time settlement point price, the index that carries the scarcity a pause is written on. \emph{A day-ahead-indexed site has the easier job, since its price is posted the afternoon before and nothing it decides on is stale, and it still finished behind in all five years on both fleets, by \$8,100 to \$55,000 per megawatt-year of total margin before grid charges. The spikes it would have paused on were never in its price, and in four years of five the day-ahead index also cost more than real time, by 1 to 9 percent of the real-time mean, a difference measured after each year closed rather than a price anyone quoted in advance. The difference is not fixed: in 2021 the day-ahead index was 2 percent cheaper and lost anyway, and in 2023 and 2024 the difference, not the spikes, decided the result. The book's rule is: if you index, index real time.}
\item[Real-Time Market (RTM)] \hypertarget{glossary:real-time-market}{}The market that balances deviations from day-ahead positions; since RTC+B it clears ancillary-service capacity and energy together. ERCOT runs SCED roughly every five minutes, repricing both as actual conditions unfold; settlement is on 15-minute intervals. \emph{Real-time prices are far more volatile than day-ahead prices, and most large flexible loads, miners included, are exposed to them.}
\item[Regional Transmission Organization (RTO)] \hypertarget{glossary:regional-transmission-organization}{}A neutral, non-profit operator that runs the grid and administers wholesale markets across a multi-state region, the form FERC Order No. 2000 (1999) encouraged. \emph{The same kind of institution as an ISO; see Independent System Operator.}
\item[Regulation (Reg-Up, Reg-Down)] \hypertarget{glossary:regulation}{}The fine motor control of the grid: ERCOT's Load Frequency Control sends a correction signal every four seconds and regulation resources must follow it within five, nudging output or consumption up (Reg-Up) or down (Reg-Down) to hold the grid at 60~Hz between economic dispatch runs. \emph{The one ancillary service a relay-only load cannot sell: only a Controllable Load Resource can.}
\item[reliability monitor] \hypertarget{glossary:reliability-monitor}{}The function that checks compliance with grid reliability rules and supports the commission's enforcement work. ERCOT performs it under the commission's direction. \emph{Distinct from the market monitor, Potomac Economics, which examines market conduct and design.}
\item[requirements contract] \hypertarget{glossary:requirements-contract}{}A retail arrangement under which a supplier delivers whatever the site consumes, priced at ERCOT's settlement point price plus a retail adder plus pass-through charges. \emph{The index base can carry a fixed-price block on top: Riot's Corsicana layered 25~MW at \$43.95/MWh for 2025 through 2027, deliberately excluding July and August.}
\item[reserve-scarcity adder] \hypertarget{glossary:reserve-scarcity-adder}{}Charges ERCOT added on top of the energy price when reserves ran short, riding on the settlement price, so a site could be billed far above the underlying energy price for the same interval. The largest of them, the ORDC adder, was retired by RTC+B; the reliability-deployment adder remains. \emph{From half past five to a quarter past eight on the morning of January 13, 2025, before the redesign, LZ\_WEST's five-minute prints stayed below \$78 while the settlement point price billed for those intervals reached \$187.77, 2.6 times the highest print in its interval. A loop deciding on bare five-minute prints would have hashed straight through: an S19j~Pro fleet ran at a loss for an hour and three-quarters that no print revealed. Chapter~3 owns the mechanism (see ORDC); Chapter~5 shows the operating consequence.}
\item[resource entity] \hypertarget{glossary:resource-entity}{}ERCOT's name for whoever owns or controls a generator, a battery or a load resource. \emph{Registration costs \$500; the third of the four gates a mine must pass to sell a megawatt-hour into ERCOT.}
\item[resource node] \hypertarget{glossary:resource-node}{}The settlement point of a generator: its own bus, paid the time-weighted average of the five-minute LMPs there. \emph{A grid battery's charging energy already settles at its resource node rather than at a zone, and an approved 2024 rule change will extend nodal settlement to every Controllable Load Resource that is not an aggregation of small sites.}
\item[Responsive Reserve Service (RRS)] \hypertarget{glossary:responsive-reserve-service}{}The grid's crash barrier: capacity that arrests a falling frequency in the first seconds after a big generator trips, deployed automatically on frequency, and also by ERCOT instruction. \emph{Its batteries and load relays act within 0.25 to 0.5 seconds and the generators' governors from about three and a half seconds (Chapter~2); it must be sustained for 15 minutes to an hour, depending on the sub-form. Three sub-forms: primary frequency response from governors, fast frequency response (almost entirely from batteries) and the under-frequency relay on loads, which is how a miner can hold this role.}
\item[retail adder] \hypertarget{glossary:retail-adder}{}A fixed dollars-per-megawatt-hour margin the retail electric provider charges on top of the index, covering its costs, its credit exposure to the buyer and its profit. \emph{A negotiated commercial term, scaled to the buyer's size and creditworthiness. Not the same thing as ERCOT's reserve-scarcity adders, which ride inside the index and are paid by everyone; since RTC+B only the reliability-deployment adder remains. Texas price-disclosure rules stop at small commercial customers, so no benchmark is published for large loads, and the one public contract in Chapter~6 redacts its Adder Charge.}
\item[retail choice] \hypertarget{glossary:retail-choice}{}The right of a consumer to pick its own retail electricity supplier rather than buy from the local monopoly. Texas opened retail choice in January 2002 under Senate Bill 7 (1999). \emph{Four ERCOT load zones belong to utilities that did not open to retail competition: Austin Energy, CPS Energy, LCRA and Rayburn.}
\item[rider] \hypertarget{glossary:rider}{}The line on a wires utility's tariff through which it passes the wholesale transmission bill down to its own customers, trued up twice a year and split among customer classes by the utility's last rate case. \emph{Because the rider is set per utility, the same megawatt of 4CP demand at transmission voltage was worth about \$40,700 to \$66,700 per megawatt-year across the four investor-owned utilities in September 2025.}
\item[SCED] \hypertarget{glossary:sced}{}Security-Constrained Economic Dispatch, ERCOT's dispatch engine. It runs roughly every five minutes, setting each node's price and telling resources what to produce. \emph{Since RTC+B each SCED run procures energy and ancillary services together, not energy alone.}
\item[security budget] \hypertarget{glossary:security-budget}{}What the whole network pays its miners in a year, which is what the network spends on its own security: the block subsidy plus transaction fees. \emph{In 2025 it was \$17.4 billion, of which \$17.2 billion was subsidy and \$0.18 billion, one percent, was fees. It is spending, not a fund anyone holds, and it will not switch from subsidy to fees on some date: it is a subsidy that halves in bitcoin terms every four years plus whatever fees the users of block space are willing to pay. Not the security of Security-Constrained Economic Dispatch, which is the grid's word for holding the network inside its operating limits.}
\item[self-generator] \hypertarget{glossary:self-generator}{}The commission registration class a generator takes if it means only to serve itself and its plant is a megawatt or more. \emph{The book uses the word for that registration class only: a plant behind the meter that serves only its host is described as part of a Private Use Network, and called a self-generator only where its registration is meant.}
\item[sell-back clause] \hypertarget{glossary:sell-back-clause}{}The clause in a fixed-price contract that lets the buyer elect not to consume its contracted power and sell it back instead, taking credits against future power bills whenever the spot price makes that the better trade. \emph{A miner who gives up resale rights signs away the asset this book is about. The sell-back and the pause are the same option at the same strike, in different paperwork: its value does not depend on which contract permits it; only the baseline moves. Riot's fixed-price contracts with sell-back rights produced \$54 million of curtailment credits in 2025, part of them the contract's own margin.}
\item[Senate Bill 6 (SB 6)] \hypertarget{glossary:senate-bill-6}{}The Texas law enacted in June 2025 in answer to a wave of proposed data centers and other large loads that had ERCOT planning for a 2030 peak of 130 to 148~GW, existing load included, against a 2024 peak of about 85~GW. Most of it is about interconnecting loads of 75~MW or more; Section 6 ordered the commission to evaluate the 4CP rule and to amend its rules by December 31, 2026. \emph{It also adds the large-load demand management service alongside ERS and requires the wires utility to have a protocol, with any equipment it needs, to curtail new transmission-connected loads during firm load shed.}
\item[Settlement Only Generator (SOG)] \hypertarget{glossary:settlement-only-generator}{}An ERCOT resource class paid for what it exports and nothing else: it may not offer energy, may not be committed or dispatched by ERCOT, and may not sell ancillary services. \emph{The interconnection process applies to a Settlement Only Generator of any size, unlike other units where the floor is one megawatt.}
\item[settlement point] \hypertarget{glossary:settlement-point}{}The place where money actually changes hands, priced from the underlying LMPs over each 15-minute settlement interval. Nobody is billed on raw bus prices. \emph{Resource nodes, load zones and hubs are three different aggregations of the same five-minute prices.}
\item[settlement point price] \hypertarget{glossary:settlement-point-price}{}The 15-minute price ERCOT actually bills a site on, built from the LMPs over the interval. ERCOT publishes a fresh five-minute price three times per settlement interval, so the loop decides on the five-minute print but the money settles here. \emph{Not the bare energy price: the reliability-deployment adder rides on top, and until RTC+B the ORDC adder did too. Not the day-ahead price either. Chapter~4 holds the real-time one at a single laboratory value so the other variables can move alone; a real ERCOT miner faces a different price every quarter-hour. See LZ\_WEST for the book's price anchor. Written in full as the real-time settlement point price where the real-time one is specifically meant, which is ERCOT's own form.}
\item[standing revenue] \hypertarget{glossary:standing-revenue}{}The book's measure of what an ancillary award is worth: what one megawatt earned by holding that award through every hour of the year, before any energy consequences. \emph{It counts capacity payments only, one product at a time. Summed across the whole product stack, from four products in 2021 to five by 2025, it fell 61-fold, from \$6.9 million per megawatt-year in 2021 to about \$113,000 in 2025. That sum measures the market, not what a megawatt could earn, since a megawatt does one thing in one quarter-hour; the best single product in 2021, RRS, paid \$2.7 million.}
\item[station power] \hypertarget{glossary:station-power}{}The electricity a power plant needs to run its own fuel pumps, control systems and cooling equipment. It normally comes from the grid itself, which is why a plant cannot produce electricity until it first receives some. \emph{The reason a black start has to begin with units that can start with no outside power.}
\item[stock power] \hypertarget{glossary:stock-power}{}The machine's factory operating point, the single power level a fleet runs at before the power target opens it into a curve. \emph{Not the same as nameplate efficiency: the measured autotuned curves sit at 28.5 and 12.8~J/TH at stock against nameplates of 29.5 and 13.5. Because the strike divides by joules per terahash, the better curves raise the pause threshold, about 3 percent for the S19j~Pro and 6 percent for the S21~XP. Chapter~9's ledger prices at the curve and its per-fleet tables at the nameplate, so the two do not add up to the dollar. A fleet on stock firmware has one operating point, not a curve.}
\item[stranded energy] \hypertarget{glossary:stranded-energy}{}Power generated where no transmission line or nearby customer can take it, so it sells for almost nothing or settles at a negative price. \emph{In ERCOT it is mostly West Texas wind and solar, built where the resource is best and far from the cities that would consume it: the energy is less wasted than priced at nothing, which is why the book's cheapest hours and its negative prices sit in the same zone. Written stranded renewable energy where the source matters; that is the same term, not a second one.}
\item[strike price (the strike)] \hypertarget{glossary:strike-price}{}The electricity price at which a machine stops paying to mine: hashprice times 41,667, divided by the model's joules per terahash. It moves with hashprice, and so must the pause threshold. By extension the book also calls the layer that pauses at this price, and what that layer earns, the strike. \emph{One strike per machine model, never one blended strike for the site. The book's backtests feed it the previous day's hashprice, because a day's network revenue is not final until the day closes; a live system can use a fresher figure. At nameplate efficiency and each year's mean hashprice the S19j~Pro's strike ran about \$440, \$172, \$104, \$90 and \$72 over the five years, and the S21~XP's from about \$960 to \$157. The strike is the breakeven price used as a live tripwire against the real-time price, not a day-ahead schedule.}
\item[summer window] \hypertarget{glossary:summer-window}{}The block of hours the site holds itself off the meter to catch the system peaks: weekdays from 4 to 6 p.m., June through September, about 175 hours a year. \emph{Chapter~8's wider rule, every day from 2 to 6 p.m. and about 490 hours a summer, is a different rule and loses money in all five years at this site. Distinct from 4CP, which is the charge the window is dodging.}
\item[supply responsibility] \hypertarget{glossary:supply-responsibility}{}Under the rules in force until December 2025, the hourly obligation a day-ahead ancillary award created: assigned to a specific resource and carried through the operating day. For a load selling Responsive Reserve the site had to be loaded, and its reserve was measured as its metered consumption over the last five minutes. \emph{This is what made the meter reading the product: the moment a committed load paused, its product ceased to exist. RTC+B replaced the physical chain with a financial position on December 5, 2025, but a relay-only load that still holds a real-time RRS or ECRS award must be loaded, and its reserve is still measured as its consumption.}
\item[System Lambda] \hypertarget{glossary:system-lambda}{}The energy component of the LMP: the cost of the next megawatt-hour anywhere on the system, before transmission limits are counted. \emph{ERCOT's own term, the one its protocols use, posted alongside every LMP. The book uses it everywhere and never writes System Energy Price. See Locational Marginal Price, congestion cost.}
\item[system-wide offer cap] \hypertarget{glossary:system-wide-offer-cap}{}The ceiling on ERCOT energy prices, \$5,000 per megawatt-hour, lowered from \$9,000 after Winter Storm Uri. \emph{In 2025 a typical ERCOT hour cleared somewhere between \$20 and \$40 per megawatt-hour, so the cap is more than a hundred times a normal price. Prices near it arrive rarely, a handful of hours in most years, but those hours are catastrophic to buy through and lucrative to avoid, and the whole economic case for miner flexibility rests on that asymmetry. Two circuit breakers sit under it: the Emergency Pricing Program, and a fall to \$2,000 for the rest of any year in which a new peaking plant's earnings above fuel pass \$315,000 per megawatt.}
\item[take-or-pay] \hypertarget{glossary:take-or-pay}{}A contract under which the buyer must accept and pay for a stated minimum volume every year whether it consumes the power or not. \emph{Cipher's Odessa PPA sets that minimum at 66.7 percent of the site's 207~MW capacity. Take-or-pay without a sell-back clause is the one contract shape that severs the machine's off switch from the market.}
\item[telemetry gateway (RTU, remote terminal unit)] \hypertarget{glossary:telemetry-gateway}{}The device at the site that reads the meters and breaker states and reports them to the QSE's control center. Substation engineers call it a remote terminal unit. \emph{It reports over a secured private line, because ERCOT requires that last mile to be a protected network rather than the open internet. Upward flow is every two seconds, to three percent accuracy or better.}
\item[transaction fees] \hypertarget{glossary:transaction-fees}{}The amounts bitcoin users attach to their own transactions, paid to the miner of the block that includes them. \emph{Through 2025 and 2026 they typically contributed about one percent of miner revenue or less. Fees are the only part of the daily pot that is not on a fixed schedule; in hashprice terms they are noise, with occasional spikes, not a trend.}
\item[transmission] \hypertarget{glossary:transmission}{}The bulk transport of electricity at high voltages over regional distances, the tall tower lines crossing the countryside. High voltages move power a long way with minimal losses. \emph{Not the same as distribution, which is the local low-voltage delivery after the step-down transformer.}
\item[transmission and distribution utility (TDU)] \hypertarget{glossary:transmission-and-distribution-utility}{}The local wires utility that bills delivery under a tariff, passed through to the miner's bill unchanged. \emph{The book's plain-English phrase for the same company is wires utility: the regulated company that owns the local lines and bills for carrying power over them.}
\item[under-frequency load shedding (UFLS)] \hypertarget{glossary:under-frequency-load-shedding}{}The automatic last resort beneath all the reserves: relays on the wires utilities' circuits that disconnect blocks of customers by themselves when frequency falls to set points, at least five percent of load at 59.3~Hz in ERCOT and more in steps below. \emph{Not the under-frequency relay of a paid load resource, set no lower than 59.70~Hz: customers shed this way neither choose nor are paid. In Uri frequency bottomed four thousandths of a hertz above the first step, and relays cut off about 276~MW even so.}
\item[under-frequency relay (UFR)] \hypertarget{glossary:under-frequency-relay}{}The bluntest way a load can sell Responsive Reserve: a hardware tripwire, armed when the site holds an RRS award, that disconnects the whole site within half a second when frequency falls to its setting, with nothing in software between the frequency and the breaker. \emph{The trip setting may be set no lower than 59.70~Hz. This is a breaker trip commanded by the grid's frequency, not the machine's own pause.}
\item[uplift] \hypertarget{glossary:uplift}{}The set of costs ERCOT spreads over every megawatt-hour of load because no market price recovers them, its own administration fee among them. \emph{The market monitor's 2025 figure, the line Chapter~6's ledger carries, is \$1.38/MWh. ERCOT allocates it in proportion to metered load, so it follows the meter rather than generation; the administration fee inside it is charged on metered load alone.}
\item[utility frequency (Hz)] \hypertarget{glossary:utility-frequency}{}The number of alternating current cycles per second on a synchronized AC grid, measured in hertz. The standard is 60~Hz in the United States and 50~Hz in many other countries, and it is the direct read-out of whether supply matches demand. \emph{Direction matters: below 60~Hz supply is short of demand, above it supply is in excess.}
\item[vertically integrated utility monopoly] \hypertarget{glossary:vertically-integrated-utility-monopoly}{}The pre-restructuring model in which one company in each region owned the power plants, the wires and the relationship with every customer, with prices set by regulators rather than by a market. State Public Utility Commissions approved the investments and granted the utility an allowed rate of return. \emph{Still the arrangement for the roughly one third of U.S. electricity consumers who live outside an ISO or RTO region.}
\item[V-floor] \hypertarget{glossary:v-floor}{}The point where a machine's power supply reaches its minimum output voltage. Below it a further power cut can only lower the clock, and the surplus voltage turns into heat, so efficiency gets worse instead of better: my reading, supported by modeling work at Braiins. \emph{The floor is a voltage limit, not a fixed wattage. It need not coincide with a firmware minimum or the best-efficiency point of a fleet average. The lowest power target the loop sets is a separate choice: see modulation floor.}
\item[Winter Storm Fern (Fern)] \hypertarget{glossary:winter-storm-fern}{}The January 2026 storm that froze two-thirds of the country and gave flexible load its strongest demonstration yet, introduced in Chapter~2 and returned to in Chapters~3, 6 and 10. ERCOT's post-storm analysis put underlying demand on the worst morning far past the all-time winter record, about 9,000~MW above what the meters actually recorded, and the grid came through with no emergency alert, no conservation appeal, and no load shed. \emph{ERCOT credited about 4,100~MW of that reduction to cryptocurrency miners, second only to oil and gas and ahead of steel, data centers and hydrogen. The network felt it too: about 200~EH/s went offline at U.S.-facing pools. See flexible load, firm load, difficulty adjustment, Winter Storm Uri.}
\item[Winter Storm Uri (Uri)] \hypertarget{glossary:winter-storm-uri}{}The February 2021 Texas storm that held real-time prices at the cap for days, the case study of Chapter~2 and the storm year of the ledger. \emph{Nearly half of the five-year value of flexibility the planned site earned falls in that one week, and seven tenths of the hindsight ceiling's, on the reserve convention that books it; 94 percent of 2021's strike value falls there too. It is also why the system-wide offer cap came down from \$9,000 to \$5,000.}
\item[Withdrawal-Limited Private Use Network (WLPUN)] \hypertarget{glossary:withdrawal-limited-private-use-network}{}A Private Use Network studied with the plant switched off and given a hard cap on what the site may draw from the grid. The site may then build up to the plant's output on top of that cap, so a large load that brings its own generation can connect before the transmission upgrades for its full draw are built. \emph{An elective route offered in ERCOT's 2026 batch study of the large-load queue, for a load that brings its own plant.}
\end{description}
\endgroup
